\section*{3 Parametric Honest Reductions}
Our first result unconditionally rules out hardness of MKtP for $\mathrm E$ under parametric honest Turing reductions. For many results on hardness of MCSP and consequences thereof, MKtP is a ``test case'' where analogous results can be shown unconditionally and where the ideas are often simpler~\cite{ref8,ref16}.
\begin{nttheorem}[Author Theorem7]
MKtP is not hard for $\mathrm E$ under parametric honest Turing reductions.
\end{nttheorem}
\begin{proof}
Let $L$ be a tally set that is in $\mathrm E$ but not in $\mathrm P$. The existence of such a tally set $L$ follows by a standard diagonalization argument.
Now suppose, for the purpose of contradiction, that MKtP is hard for $\mathrm E$ under parametric honest Turing reductions. This implies that there is an oracle Turing machine $M$ running in time at most $cn^k$, where $c$ and $k\geq1$ are constants, such that $M$ decides $L$ correctly with oracle access to MKtP, and moreover there is a constant $\epsilon>0$ such that each query made by $M$ on an input of length $n$ has parameter at least $n^\epsilon$. We show how to use $M$ to decide $L\in\mathrm P$, contrary to the assumption on $L$.
Consider the following polynomial-time machine $A$ that attempts to decide $L$. Given input $0^n$, $A$ runs the oracle machine $M$, answering each query made by $M$ with ``yes''. When the simulation of $M$ is complete, $A$ accepts iff $M$ accepts. Since $M$ is polynomial-time, so is $A$. We show that $A$ correctly decides $L$ for large enough $n$.
Given input $0^n$, $M$ makes some sequence of oracle queries $q_1,q_2\ldots q_t$ during the simulation by $A$, where $0\leq t\leq cn^k$. We show that for each $i$, $q_i$ has $Kt$ complexity $O(\log(n))$. Observe that $q_i$ can be generated by the universal machine $U$ implicit in the definition of $Kt$ complexity in time $\poly(n)$ given descriptions of $i,n$ and the oracle machine $M$: it simply needs to run $M$ on $0^n$ answering the first $i-1$ queries made by $M$ with ``yes'', and then output the $i$'th query. The description length required is $O(\log(n))$ since $i\leq cn^k$, and moreover the time taken by $U$ is polynomial in $n$, hence using the definition of $Kt$ complexity, we have an $O(\log(n))$ upper bound on the $Kt$ complexity of $q_i$. Let $C$ be a constant such that the $Kt$ complexity of each $q_i$ is bounded above by $C\log(n)$.
By the assumption that the reduction is parametric honest, the parameter of $q_i$ is at least $n^\epsilon$, which for large enough $n$ is greater than $C\log(n)$. Hence for large enough $n$ and for each $i$, it follows that $A$'s assumed answer for the oracle query $q_i$ is correct for each $i$, and hence that $A$ outputs the same answer as $M$ would given the true oracle MKtP. Since $M$ is an oracle Turing machine correctly implementing a reduction from $L$ to MKtP, $A$ is a polynomial-time Turing machine correctly deciding $L$ for large enough $n$, contradicting the assumption on $L$.
\end{proof}
We now re-state and prove Theorem1.
\begin{nttheorem}[Author Theorem8]
If MCSP is NP-hard under parametric honest Turing reductions, then $\mathrm E\not\subseteq\SIZE(\poly)$.
\end{nttheorem}
\begin{proof}
Suppose that MCSP is NP-hard under parametric honest Turing reductions. Let $M$ be a polynomial-time oracle Turing reduction implementing a parametric honest reduction from SAT to MCSP, and let $\epsilon$ be a constant such that the parameter of any query made by $M$ on any SAT instance of length $x$ is bounded below by $|x|^\epsilon$.
Consider the following polynomial-time machine $A$ which attempts to decide SAT. On an input $x$ of SAT, $A$ runs the oracle Turing machine $M$. Assume wlog that all queries asked by $M$ have length a power of two--any queries for which this is not the case may be eliminated by assuming the answer is `no'. For every query asked by $M$ with length a power of two, $A$ assumes the answer is ``yes'', and continues the simulation of $M$. When the simulation concludes, $A$ accepts iff $M$ accepts.
Clearly $A$ runs in polynomial time. If $A$ decides SAT correctly, then we have that $\mathrm{NP}=\mathrm P$. Hence the Polynomial Hierarchy collapses to $\mathrm P$. But this implies that $\mathrm E\not\subseteq\SIZE(\poly)$; if the contrary were true, then $\mathrm E$ would collapse to the Polynomial Hierarchy using the Karp--Lipton theorem for $\mathrm E$, and hence to $\mathrm P$, contradicting the hierarchy theorem for deterministic time.
Suppose now that $\mathrm{NP}\ne\mathrm P$. In particular, this means that there are infinitely many instances on which $A$ does not decide SAT correctly. If not, we could hardcode the finitely many instances on which $A$ fails to compute SAT into a new polynomial-time machine $A'$ which runs in polynomial time and solves SAT correctly on all instances, contradicting the assumption that $\mathrm{NP}\ne\mathrm P$.
Now, we can use the main result of Gutfreund, Shaltiel and Ta-Shma. They show that if $\mathrm{NP}\ne\mathrm P$, then for any polynomial-time algorithm $A$ that fails to compute SAT, there is a polynomial-time Turing machine $B$ and a constant $d$, such on input $1^n$, $B$ outputs three instances $x_1,x_2,x_3$ each of length $n$ or length $n^d$ such that for infinitely many $n$, $A$ fails to decide SAT correctly on at least one of these three instances.
We show how to use $B$ to argue again that $\mathrm E\not\subseteq\SIZE(\poly)$. Consider the three instances $x_1,x_2,x_3$ output by $B$ on input $1^n$. For each $i$, $1\leq i\leq3$, let $Q_i$ be the concatenation of all queries made by $M$ when $M$ is simulated by $A$ on input $x_i$. Let $Y$ be the concatenation of $Q_1,Q_2,Q_3$. Note that $|Y|=\poly(n)$, since $M$ is a polynomial-time oracle Turing machine. Also note that each query made by $M$ on an input of length $n$ has parameter at least $n^\epsilon$, by the parametric honesty of the reduction, as long as all previous queries were answered correctly. Each such query was assumed by $A$ to have a ``yes'' answer. It could not be the case that all such answers were correct for all queries made by $M$ on inputs $x_1,x_2,x_3$--if they were, then $A$ would correctly solve SAT on $x_1,x_2,x_3$, using the fact that $M$ implements a correct polynomial-time Turing reduction from SAT.
Thus we get that there is a first query made by $M$ on one of $x_1,x_2,x_3$ that has parameter at least $n^\epsilon$ and is wrongly answered ``yes''. Fix such a query $y$. Since $M$ implements a Turing reduction from SAT to MCSP, $y$ is the truth table of a Boolean function on $O(\log(n))$ bits that requires circuits of size $n^\epsilon$. We have that $y$ is a substring of $Y$, hence $Y$ must also require circuits of size $\Omega(n^\epsilon)$. We are using here the simple fact that if $y$ is a substring (with length a power of two) of the truth table $Y$ of a Boolean function with circuits of size $s$, then $y$ is the truth table of a Boolean function with circuits of size $O(s)$.
Hence, for infinitely many $n$, on input $1^n$, we can compute the truth table $Y$ of a Boolean function on $O(\log(n))$ bits such that the function requires circuits of size $n^\epsilon$, i.e., of exponental size in the length of the input to the Boolean function. This implies $\mathrm E\not\subseteq\SIZE(\poly)$.
\end{proof}
